
\documentclass[10pt,aspectratio=169,mathserif]{beamer}		
%设置为 Beamer 文档类型，设置字体为 10pt，长宽比为16:9，数学字体为 serif 风格
\usepackage[utf8]{inputenc}
\usepackage{ctex}
\usepackage{lmodern}      % 字体支持
\usepackage{graphicx}     % 图片支持
\usepackage{amsmath}      % 数学公式
\usepackage{tikz}         % 绘图包（必须添加！）
\usetheme{Madrid}         % 主题
\usepackage[T1]{fontenc}        % 更好的字体编码
\usepackage{lmodern}            % Latin Modern 字体，包含完整的斜体和粗体
\usepackage{amsmath,amssymb}    % 数学符号
\usepackage{xeCJK}              % 如果你需要中文字体支持
% 添加这个包在导言区
\usepackage{float}
\usepackage{booktabs} 
\usetikzlibrary{shapes, arrows.meta, positioning}
%%%%-----导入宏包-----%%%%
\usepackage{whpu}			%导入 whpu 模板宏包
\usepackage{ctex}			%导入 ctex 宏包，添加中文支持
\usepackage{amsmath,amsfonts,amssymb,bm}   %导入数学公式所需宏包
\usepackage{color}			 %字体颜色支持
\usepackage{graphicx,hyperref,url}
\usepackage{metalogo}	% 非必须
%% 上文引用的包可按实际情况自行增删
%%%%%%%%%%%%%%%%%%
\usepackage{fontspec}
\usepackage{xeCJK}
% \setCJKmainfont{Source Han Sans SC}

% 直接在文档中嵌入BibTeX条目，避免外部文件问题


\beamertemplateballitem		%设置 Beamer 主题

%%%%------------------------%%%%%
\catcode`\。=\active         %或者=13
\newcommand{。}{．}				
%将正文中的“。”号转换为“.”。中文标点国家规范建议科技文献中的句号用圆点替代
%%%%%%%%%%%%%%%%%%%%%

%%%%----首页信息设置----%%%%
\title[武汉轻工大学]{系数求解与预测优化}
\subtitle{——课程设计答辩}			
%%%%----标题设置


\author[你的名字]{
  韦一敏 \\\medskip }
%%%%----个人信息设置
  
\institute[IOPP]{
  数学与计算机学院 \\ 
  武汉轻工大学}
%%%%----机构信息

\date[2025.12]{
  2025年12月}
%%%%----日期信息
  
\begin{document}

\begin{frame}
	\titlepage
\end{frame}				%生成标题页

\section{目录}
\begin{frame}
	\frametitle{目录}
	\tableofcontents
\end{frame}			

\section{背景介绍}
\begin{frame}
    \frametitle{项目背景与目标}
    通过数据预处理、模型训练与评估，旨在为招生决策提供量化工具，同时为算法选择提供实证参考。
\end{frame}

\section{流程图}
\begin{frame}{简洁横向流程图}
\centering
\begin{tikzpicture}[
    node distance=2.2cm,
    box/.style={
        draw,
        rectangle,
        rounded corners=4pt,
        minimum width=2.2cm,
        minimum height=1cm,
        align=center,
        fill=blue!10
    },
    arrow/.style={->, >=stealth, thick}
]

% 水平排列
\node[box] (1) {步骤1};
\node[box, right of=1] (2) {步骤2};
\node[box, right of=2] (3) {步骤3};
\node[box, right of=3] (4) {步骤4};

% 连接
\draw[arrow] (1) -- (2);
\draw[arrow] (2) -- (3);
\draw[arrow] (3) -- (4);

% 标签
\node[above of=1, yshift=-0.5cm] {输入};
\node[above of=4, yshift=-0.5cm] {输出};

\end{tikzpicture}

\vspace{1cm}
\begin{columns}
\column{0.5\textwidth}
\textbf{优点：}
\begin{itemize}
    \item 结构清晰
    \item 易于理解
    \item 节省空间
\end{itemize}

\column{0.5\textwidth}
\textbf{应用：}
\begin{itemize}
    \item 算法流程
    \item 工作流程
    \item 系统架构
\end{itemize}
\end{columns}

\end{frame}




% 目录页（可选）
\begin{frame}{目录}
  \tableofcontents
\end{frame}

% 内容部分
\section{第一节标题}
\begin{frame}{第一张幻灯片标题}
  你的内容...

  我不写了
\end{frame}

\section{第二节标题}
\begin{frame}{第二张幻灯片标题}
  更多内容...

  你自己AI
\end{frame}

\section{第三节标题}
\begin{frame}{第三张幻灯片标题}
 
    好困

\end{frame}

\section{第四节标题}
\begin{frame}{第四张幻灯片标题}
 
    好好学习

    天天向上

\end{frame}


\begin{frame}
    \frametitle{实验结果总结与解释}
    
    而后钱钱钱钱钱钱钱
 \end{frame}   
 
 % 参考文献页（如果需要）
% 只在需要时才取消注释下面几行
%\begin{frame}{参考文献}
%\bibliographystyle{plain}
%\bibliography{references}  % 确保有这个文件
%\end{frame}

 \begin{frame}
    \centering
    \definecolor{whpublue}{RGB}{0, 86, 166}
    \vspace*{\fill}
    
    {\fontsize{36}{40}\selectfont\itshape\bfseries\textcolor{whpublue}{THANK YOU}}
    
    \vspace*{\fill}
    
\end{frame}

\end{document}